\documentclass[11pt,a4paper]{article}
\usepackage[margin=1.8cm,headheight=48pt,headsep=12pt,footskip=22pt]{geometry}
\usepackage{fontspec}
\setmainfont{Noto Serif}
\setsansfont{Noto Sans}
\setmonofont{DejaVu Sans Mono}
\usepackage{amsmath,amssymb,mathtools}
\usepackage{enumitem,tabularx,booktabs,array,multicol}
\usepackage{xcolor,graphicx,tikz,fancyhdr,lastpage}
\usepackage{listings,tcolorbox,hyperref}
\hypersetup{colorlinks=true,linkcolor=blue!60!black,urlcolor=blue!60!black}
\definecolor{hcmusblue}{RGB}{28,55,111}
\definecolor{lightblue}{RGB}{237,242,250}
\definecolor{codebg}{RGB}{247,248,250}
\setlength{\parindent}{0pt}
\setlength{\parskip}{4pt}
\setlist[enumerate]{leftmargin=*,itemsep=5pt,topsep=4pt}
\setlist[itemize]{leftmargin=*,itemsep=2pt,topsep=2pt}
\lstset{language=C++,basicstyle=\ttfamily\footnotesize,backgroundcolor=\color{codebg},frame=single,rulecolor=\color{black!15},breaklines=true,showstringspaces=false,columns=fullflexible,keywordstyle=\color{hcmusblue}\bfseries,commentstyle=\color{green!40!black},numbers=left,numberstyle=\tiny\color{black!45},xleftmargin=1.5em,framexleftmargin=1.2em}
\pagestyle{fancy}
\fancyhf{}
\fancyhead[L]{\includegraphics[height=1.05cm,keepaspectratio]{assets/images/institution_logo.png}\hspace{0.45em}\begin{minipage}[c]{0.64\textwidth}\textbf{TRƯỜNG ĐẠI HỌC KHOA HỌC TỰ NHIÊN}\\[-1pt]\small ĐHQG TP. HỒ CHÍ MINH\end{minipage}}
\fancyhead[R]{\small\textcolor{hcmusblue}{TÀI LIỆU LUYỆN TẬP}}
\fancyfoot[L]{\scriptsize Biên soạn độc lập - Không phải đề thi chính thức}
\fancyfoot[R]{\scriptsize Trang \thepage}
\renewcommand{\headrulewidth}{0.5pt}
\renewcommand{\footrulewidth}{0.3pt}
\newcommand{\examtitle}[4]{\begin{center}{\Large\bfseries\color{hcmusblue}#1}\par\vspace{3pt}{\large\bfseries #2}\par\vspace{3pt}\small #3\quad\textbullet\quad #4\end{center}\vspace{5pt}\hrule\vspace{8pt}}
\newcommand{\studentline}{\begin{tabularx}{\textwidth}{X X}\textbf{Họ và tên:} \hrulefill & \textbf{MSSV:} \hrulefill\\[7pt]\textbf{Lớp:} \hrulefill & \textbf{Ngày làm:} \hrulefill\end{tabularx}\vspace{8pt}}
\newcommand{\points}[1]{\hfill\textbf{(#1 điểm)}}
\newtcolorbox{notebox}{colback=lightblue,colframe=hcmusblue!55,boxrule=0.5pt,arc=1.5mm,left=2mm,right=2mm,top=1mm,bottom=1mm}
\begin{document}
\examtitle{ĐỀ LUYỆN TẬP 05 - CƠ SỞ LẬP TRÌNH}{CƠ SỞ LẬP TRÌNH - HCMUS (theo lịch học tuần 1--5)}{90 phút - 10 điểm - C++11 - Không sử dụng tài liệu}{Giữa kỳ 1 - HK1 2026--2027}

\studentline
\begin{notebox}\textbf{Phạm vi:} Vòng lặp \texttt{while}, đệ quy, tham chiếu, hàm xử lý mảng 1D và kiểm tra điều kiện trên ma trận 2D.\end{notebox}
\begin{enumerate}
\item Viết hàm đảo ngược chữ số của số nguyên không âm $n$ bằng vòng lặp \texttt{while}. Giả sử kết quả không vượt giới hạn kiểu \texttt{int}. Ví dụ $n=2048$ cho kết quả 8402. \points{2}
\item Viết hàm đệ quy \texttt{factorial(n)} tính $n!$ với $0\le n\le12$. Nêu điều kiện dừng. \points{2}
\item Viết hàm nhận mảng số nguyên có ít nhất một phần tử và tìm giá trị nhỏ nhất, lớn nhất. Trả cả hai kết quả thông qua tham số tham chiếu \texttt{mn}, \texttt{mx}. \points{2}
Prototype: \texttt{void findMinMax(const int a[], int n, int \&mn, int \&mx);}
\item Viết hàm chèn giá trị \texttt{x} vào vị trí \texttt{pos} của mảng có sức chứa tối đa \texttt{capacity}. Vị trí hợp lệ từ 0 đến $n$. Nếu mảng đầy hoặc vị trí không hợp lệ, không thay đổi mảng và trả về \texttt{false}; ngược lại dịch phần tử sang phải, cập nhật $n$ và trả về \texttt{true}. \points{2}
Prototype: \texttt{bool insertAt(int a[], int \&n, int capacity, int pos, int x);}
\item Viết hàm đếm số cột mà tất cả phần tử trong cột đó đều không âm trong ma trận $R\times C$. Giả sử $R,C>0$. \points{2}
Prototype: \texttt{int countNonNegativeColumns(const int a[][MAX], int R, int C);}
\end{enumerate}
\vfill\begin{center}\textit{--- Hết ---}\end{center}

\end{document}
